\section{结果分析(系统效果与效率)}

\begin{plainbox}
\textbf{说明}:本章只写三类可复核证据:历史固定评测(检索 recall/MRR 与延迟)、当前工程运行态(测试、全量 embedding、OpenCode/MCP 闭环)和当前自动代理基线(推荐 proxy、趋势回测)。历史评测、工程 PASS 与外部质量证明分开陈述;若与真实业务场景不一致,必须以后续人工 Gold、用户试点或场景化验证补充,避免确定性外推。
\end{plainbox}

\begin{reportbox}{核心指标一览}
\begin{tableblock}
\begin{tabularx}{0.96\textwidth}{>{\bfseries}p{32mm}X}
\toprule
检索质量(历史固定评测) & 标题自检索 recall@10=\textbf{0.985},MRR@10=\textbf{0.9583},平均延迟 \textbf{141.6 ms};关键词自检索 recall@10=0.745,MRR@10=0.7045。 \\
跨语言相关性(历史定向样例) & 15 条固定中文主题查询上的 relevance@5=\textbf{1.0}(15/15),优化前为 0.667; 该结果只代表定向样例,不是全局质量证明。 \\
趋势描述边界(当前回测) & Pearson=\textbf{0.9756},但 MAE 与 sMAPE 在 2/2 个滚动切分上均未超过 naive 基线,故当前只可称\textbf{描述性趋势}。 \\
论文推荐(四路自动多基线) & `current` 100 seeds/500 recs: same-field \textbf{0.7180},shared-keyword \textbf{0.6280},semantic \textbf{0.9263},diversity \textbf{0.0823},novelty proxy \textbf{0.5885}; 相对 `dense` 的 0.6940/0.5080/0.9332/0.0682/0.5810 是工程权衡,不是全面胜出。 \\
工程运行态(2026-08-10) & 后端回归 \textbf{524 passed, 1 skipped}; 远端全量 embedding 运行态已验证 \texttt{chunk\_embeddings=367{,}861/367{,}861}、\texttt{paper\_embeddings=159{,}164/159{,}164}; OpenCode 顺序调用核查 tool + disputes resource 已通过。 \\
\bottomrule
\end{tabularx}
\end{tableblock}
\end{reportbox}

\begin{figure}[!htbp]
\centering
\figasset[0.94\textwidth]{eval_metric_dashboard.png}
\caption{核心评测指标图形化总览。数据来源:\texttt{output/eval/eval\_report.json}、\texttt{output/eval/trends\_backtest.json}、\texttt{output/eval/recommend\_bakeoff.json} 与当前验收记录;中文相关性为固定主题查询评测口径,测试与 OpenCode 为 2026-08-10 工程基线。}
\label{fig:eval-dashboard}
\end{figure}

\begin{figure}[!htbp]
\centering
\figasset[0.94\textwidth]{asset_funnel.png}
\caption{数据到智能体的资产漏斗。数据来源:\texttt{data/analysis/summary.json}、\texttt{data/processed/*.summary.json} 与 \texttt{output/eval/eval\_report.json};图中为文件资产口径,运行时数据库/向量表口径在本章表格中另列。}
\label{fig:asset-funnel}
\end{figure}

图~\ref{fig:eval-dashboard} 将核心效果指标集中展示,图~\ref{fig:asset-funnel} 则说明这些指标背后的文件资产链路。二者分别对应"系统效果"与"可复现资产"两个评审入口。

\subsection{检索质量与效率(核心测试)}
采用\textbf{自检索协议}(以论文自身标题/关键词为查询,检验能否检回该文)评测混合检索。评测结果固化于 \texttt{output/eval/eval\_report.json};当前运行时库表包含 159{,}164 篇论文和 367{,}861 个片段,该章中的 recall/MRR 仍引用 2026-06-22 历史固定评测而非 2026-08-10 重跑结果:

\begin{tableblock}
\begin{tabularx}{0.94\textwidth}{Xcccc}
\toprule
\textbf{查询方式} & \textbf{recall@1} & \textbf{recall@10} & \textbf{MRR@10} & \textbf{平均延迟} \\
\midrule
按标题(字面 + 语义臂) & 0.94 & \textbf{0.985} & 0.9583 & 141.6\,ms \\
按关键词 & 0.6913 & 0.745 & 0.7045 & 84.2\,ms \\
\bottomrule
\end{tabularx}
\end{tableblock}

标题查询 recall@10 达 \textbf{0.985}(目标论文几乎总在前 10),平均延迟 141.6 ms,体现混合检索"\textbf{效果与效率}"兼顾。关键词查询由于同义词、缩写和领域词复用更强,指标低于标题查询,但仍保持 recall@10=0.745 和 MRR@10=0.7045;本报告保留该差异,不把关键词检索夸大成标题级精确命中。

\subsection{跨语言检索相关性优化(历史定向样例)}
针对中文短查询,初版多语言嵌入存在"按语言而非主题召回"的偏差
(如查"扩散模型"却召回到无关的中文医学论文)。我们用 15 个主题查询构建
\texttt{relevance@5} 指标(top-5 是否含主题相关论文,\texttt{make} 可复现),针对性实施了一组优化:

\begin{tableblock}
\begin{tabularx}{0.94\textwidth}{lX}
\toprule
\textbf{优化} & \textbf{作用} \\
\midrule
中英术语映射 & 中文查询自动替换/补充英文等价词,把召回锚定到(以英文为主的)语料主题 \\
混合字面臂 & \texttt{websearch} 精确短语优先 + OR 兜底,兼顾精确命中与多词召回 \\
语义跨语言答案校验 & 用 e5 余弦度量"答案$\leftrightarrow$证据"接地度,中文答案据英文文献不再被误判低置信 \\
"最新"意图识别 & 识别"最新/近期"等意图,在相关结果中优先近年论文 \\
多轮对话记忆 & 指代型追问("它的局限呢")自动承接上文主题再检索 \\
\bottomrule
\end{tabularx}
\end{tableblock}

在这组固定样例上,\texttt{relevance@5} 从 \textbf{0.667 提升至 1.0}(15/15 全部命中主题),
"扩散模型""强化学习""量子计算""推荐系统"等此前跑偏的中文查询被修复。该结果是\textbf{定向样例工程优化证据},不能替代更大规模人工 qrels 或真实用户评测。

\subsection{全量向量运行态}
2026-08-10 在 2080 Ti 远端运行态中,全量 \texttt{chunk\_embeddings} 与
\texttt{paper\_embeddings} 已分别达到 367{,}861/367{,}861 与
159{,}164/159{,}164,从而关闭了"推荐缺 embedding"这一工程阻塞。这里仅使用能够由远端数据库
计数和运行日志复核的总量,不再保留缺少独立证据产物的领域覆盖消融数字。该证据证明的是
\textbf{全量资产可运行};推荐质量与用户价值仍需单独评测。

\subsection{差异化能力对比}
\begin{figure}[!htbp]
\centering
\figasset[0.88\textwidth]{claim_grounding_flow.png}
\caption{论断接地流程。数据来源:工具契约与黄金会话;输出为支持等级和可追溯证据。}
\label{fig:claim-grounding-flow}
\end{figure}

图~\ref{fig:claim-grounding-flow} 对应差异化矩阵中的"据实接地/可验证"能力。系统给出支持等级与出处,但不把语义相似度包装为形式逻辑证明。

\begin{tableblock}
\begin{tabularx}{0.96\textwidth}{>{\bfseries}p{34mm}XXXX}
\toprule
能力维度 & \textbf{SciScope} & 通用大模型 & 商用文献工具 & 普通检索 \\
\midrule
据实接地 / 可验证 & \textbf{论断核查 + 引用证据} & 易幻觉 & 部分支持 & 不支持 \\
跨语言中→英检索 & \textbf{固定中文主题样例已验证} & 依模型能力波动 & 通常较弱 & 不支持 \\
本地部署 / 数据主权 & \textbf{数据、索引、模型本地} & 查询上传第三方 & 查询上传第三方 & 仅检索外部库 \\
自主多步编排 & \textbf{可观测规划/反思 + 专业技能 + 10 类运行时能力} & 非工具化问答 & 固定流程为主 & 不支持 \\
可复现 / IP 清晰 & \textbf{\texttt{make} 重建资产} & 黑盒调用 & 平台黑盒 & 无模型资产 \\
\bottomrule
\end{tabularx}
\end{tableblock}
差异化不来自"又接了一个大模型",而来自\textbf{本地语料资产、可复现模型文件、证据接地检验和自主工具编排}共同形成的科研智能体闭环。

\subsection{趋势回测与推荐评估(评测证据包)}
历史自检索评测由 \texttt{make eval-all} 保留;当前趋势回测由
\texttt{python -m evaluation.eval\_trends\_backtest} 重建,推荐四路对照由
\texttt{python -m evaluation.eval\_recommend\_bakeoff} 在全量远端数据库重建。对应产物分别是
\texttt{output/eval/eval\_report.\{json,md\}}、\texttt{trends\_backtest.\{json,md\}} 与
\texttt{recommend\_bakeoff.\{json,md\}};不再声称单个旧入口覆盖全部当前指标。

\textbf{趋势回测}(拟合 2022--2024 → 预测 2025,1{,}998 个关键词):预测与真实的 Pearson 相关 \textbf{0.9756},但在两个固定滚动切分上,趋势模型的 MAE 与 sMAPE 均未超过 last-value / moving-average naive 基线,因此当前输出被明确降级为 \textbf{descriptive\_only=true}。它可以用于描述历史方向、波动和 burst,但\textbf{不得}在答辩或报告中写成"预测能力"或"研究方向预测结论"。

\textbf{推荐自动多基线评测}(2026-08-10,100 篇 seeds,top-5):四路方法 `popularity\_proxy / keyword\_tfidf / dense / current` 已统一落表。`current` 的 coverage 为 \textbf{1.0},same-field \textbf{0.7180},shared-keyword \textbf{0.6280},平均语义相似 \textbf{0.9263},diversity \textbf{0.0823},novelty proxy \textbf{0.5885}; 相比 pure `dense` 的 0.6940 / 0.5080 / 0.9332 / 0.0682 / 0.5810,它在同领域率、共享关键词、多样性和新颖度代理上更高,但语义相似度更低,且延迟显著更高(本次 `current` p50 约 100.6\,ms, `dense` p50 约 4.2\,ms)。这些都是\textbf{自动代理指标},只能说明推荐器之间的工程权衡,不能替代研究者满意度、用户盲评或真实科研提效证明。

\subsection{数据治理成效}
\begin{tableblock}
\begin{tabularx}{0.94\textwidth}{p{34mm}p{48mm}X}
\toprule
\textbf{治理项} & \textbf{处理} & \textbf{效果} \\
\midrule
标题优先去重 + front-matter 剔除 & 分析层与处理后语料分层治理 & 残留长标题重复组归零 \\
PDF 垃圾正文过滤 & 清 83 篇 / 46 片段 & 显著减少回答混入 PDF 乱码的风险 \\
关键词噪声过滤 & arXiv 码/泛标签/期刊名 & 趋势 top 回归真实主题 \\
摘要覆盖治理 & 140{,}183 篇有摘要 & 摘要覆盖率 83.22\%,支撑主题和检索解释 \\
开放正文补充证据 & 分析层正文片段覆盖集中于开放全文来源 & 分析层正文片段覆盖率 8.99\%; 处理后语料中 full-text 非空为 13{,}400/159{,}164(8.42\%),只用于深读补充证据 \\
\bottomrule
\end{tabularx}
\end{tableblock}
关键词清洗后,趋势热点回归真实研究主题:\textbf{retrieval augmented generation}、\textbf{large language model} 位居新兴前列;主题模型(40 主题)清晰区分 LLM、RAG、图神经网络、锂电池、蛋白结构、催化等方向。

\subsection{模型资产与工程质量}
\begin{tableblock}
\begin{tabularx}{0.94\textwidth}{Xc}
\toprule
\textbf{模型文件 / 索引 / 质量项} & \textbf{规模} \\
\midrule
片段向量 \texttt{chunk\_embeddings}(pgvector) & 367{,}861（2026-08-10 远端全量运行态） \\
论文级向量 \texttt{paper\_embeddings} & 159{,}164（2026-08-10 远端全量运行态） \\
主题模型(LDA + NMF) & 40 主题 / 模型 \\
知识图谱(作者/关键词/主题) & 客户端友好 JSON \\
后端自动化测试 & \texttt{make test-backend} 实测 524 passed、1 skipped（2026-08-10）; \texttt{make agent-smoke} 仅作历史黑盒记录,不替代当前工程基线 \\
\bottomrule
\end{tabularx}
\end{tableblock}
\begin{plainbox}
\textbf{产品化呈现}:智能体运行时采用可观测状态图编排,向终端客户端输出科研工作流阶段、耗时、会话与重试信息;TUI 据此渲染阶段条、分组时间线、会话保存与错误恢复。最终回答采用"研究结论 + 证据工具 + 过程入口"结构,工具日志不混入主聊天。2026-08-10 已新增两条工程证据:一是全量 embedding/ANN 运行态关闭了推荐阻塞,二是真实 OpenCode 已能顺序完成核查 tool 与争议 resource 读取。全部模型文件仍可由 \texttt{make agent-build} 一键重建,知识产权清晰、评委可复现。
\end{plainbox}

\subsection{开源治理、分发与托管体验}
除本地完整流水线外,SciScope 已补齐面向真实用户的交付链路。项目按 Apache License 2.0 开源,仓库提供贡献指南、安全策略、行为准则、Issue/PR 模板与 GitHub Discussions,使后续数据治理、后端、TUI、打包发布和文档优化可以以开源协作方式持续迭代。

\begin{tableblock}
\begin{tabularx}{0.96\textwidth}{>{\bfseries}p{30mm}X}
\toprule
\textbf{交付面} & \textbf{当前状态与边界} \\
\midrule
终端客户端分发 & 提供 \texttt{npm}、Homebrew 与 Scoop 三条安装路径;其中 Scoop 是当前 Windows 稳定路径,Homebrew 面向 macOS 用户,\texttt{npm} 面向已有 Node.js 的跨平台开发者。 \\
公开托管后端 & 公开体验入口采用 Render Web Service + Supabase PostgreSQL/pgvector + DeepSeek API,用于安装后快速试用和演示视频录制。 \\
能力边界 & 托管后端定位为 demo-lite 子集,受免费数据库容量与冷启动限制;完整语义覆盖、全量向量表和长期稳定服务仍以本地/自建 PostgreSQL + pgvector 部署为准。 \\
提交边界 & 最终提交按"报告与附件逐文件上传"组织,两份 PDF 独立上传;附件保留源代码、运行说明、评测/图谱/模型小型结果,不随附大型 \texttt{data/}、\texttt{models/} 与数据库运行表。 \\
\bottomrule
\end{tabularx}
\end{tableblock}

因此,本项目的产品化状态不是只停留在离线脚本:评委可通过报告和附件审阅完整方法,通过本地命令重建全量能力,也可通过已分发的 TUI 连接公开后端进行轻量体验。三者分别服务"可审阅、可复现、可试用"三个目标。
